% ================= APPENDIX: REQUIREMENTS VERIFICATION MATRIX =================
% \appendix is declared in main.tex, before the first appendix input.
\section{Requirements Verification Matrix}
\label{app:verification}

This appendix is the working record behind Table~\ref{tab:requirements}: for
each requirement, the procedure that verifies it, the criterion that closes it,
the evidence to date and the current status. Section~\ref{sec:requirements}
states \emph{what} must be true; this appendix states \emph{how it is shown},
and is updated as testing proceeds rather than written once.

Status values are: \textbf{Verified}, evidence exists and meets the criterion;
\textbf{Partial}, evidence exists on a representative article but not on the
delivered one; \textbf{Open}, no evidence yet. A requirement verified on the
1-DoF rig is \emph{Partial} with respect to the cube, because the rig is a
different article --- the distinction is kept deliberately, since collapsing it
would overstate what the integrated machine has been shown to do.

\begin{footnotesize}
\setlength{\tabcolsep}{3.5pt}
\begin{longtable}{@{} l c L{5.0cm} L{4.2cm} L{2.3cm} @{}}
\caption[Requirements verification matrix]{Requirements verification matrix.
Methods: A analysis, I inspection, T test, D demonstration.}
\label{tab:verification}\\
\toprule
ID & V & Procedure & Acceptance criterion & Status \\
\midrule
\endfirsthead
\multicolumn{5}{c}{\footnotesize\itshape \tablename~\thetable\ (continued)}\\
\toprule
ID & V & Procedure & Acceptance criterion & Status \\
\midrule
\endhead
\midrule
\multicolumn{5}{r}{\footnotesize\itshape Continued on next page}\\
\endfoot
\bottomrule
\endlastfoot

\multicolumn{5}{@{}l}{\textbf{Functional}}\\
\midrule
FUN-01 & T & Release the cube at the edge equilibrium; run the balancing controller with no operator contact. & Balance sustained; no external support at any point. & Open \\
FUN-02 & T & As FUN-01, at the corner equilibrium, with the three-axis gain set. & Balance sustained on a point contact. & Open \\
FUN-03 & T & Command a slew between two defined attitudes from a stable balance. & Target attitude reached and balance re-captured without a fall. & Open \\
FUN-04 & T & Apply a repeatable impulse to the body during balance. & Return to the equilibrium without a fall. & Partial (rig) \\
FUN-05 & T & Compare the estimator output against a known reference attitude, statically and while moving. & Estimate tracks the reference within the band recorded in Sec.~\ref{sec:test_1dof}; no external reference used in flight. & Open \\
FUN-06 & T & Balance until a wheel reaches the desaturation threshold; log the commanded bias torque. & Desaturation torque commanded; wheel speed bounded without loss of balance. & Open \\
FUN-07 & D & Exercise every transition of Table~\ref{tab:sw_modes}, including at each threshold boundary. & All transitions taken as specified; no chattering at any boundary. & Open \\
FUN-08 & D & Issue a disarm command from each operating mode in turn. & Drivers de-energised from every mode without exception. & Open \\
FUN-09 & D & Induce a fall from balance and observe the response. & Safe state entered with wheels commanded down; no continued actuation. & Open \\
FUN-10 & T & Spin up to the design speed from a face rest; actuate the brake; attempt capture. & Cube reaches an edge or corner and captures into balance. & Open \\
FUN-11 & T & Execute consecutive jump-ups with re-stabilisation between each. & Two or more consecutive hops with balance re-established after each. & Open \\
FUN-12 & D & Run the control loop with telemetry enabled and disabled; compare cycle timing. & No change in worst-case cycle time attributable to telemetry. & Open \\
\midrule

\multicolumn{5}{@{}l}{\textbf{Performance}}\\
\midrule
PER-01 & T & Initialise at rest at the upright equilibrium; release; time the balance. Run on the 1-DoF rig, then on the integrated cube at the edge equilibrium. & Duration $\geq$ \qty{60}{\second}; no evident wheel-speed drift, on each article. & Partial (rig verified; cube open) \\
PER-02 & T & As PER-01, on the integrated cube at the corner equilibrium. & Duration $\geq$ \qty{60}{\second}; no evident drift. & Open \\
PER-03 & T & Displace the rig by hand to a measured angle from the support; release. & Recovery from $\geq$ \qty{12}{\degree}, then PER-01 satisfied. & \textbf{Verified} (to \qty{24}{\degree}) \\
PER-04 & T & As PER-03, on the cube, in each of the principal tilt directions. & Recovery from $\geq$ \qty{5}{\degree}, then PER-01 or PER-02 satisfied. & Open \\
PER-05 & T & Instrument the control task; log per-cycle execution time over a full balancing run. & Loop rate \qty{1}{\kilo\hertz}; worst-case cycle $\leq$ \qty{1}{\milli\second}, not merely the mean. & Open \\
PER-06 & A, I & Evaluate wheel inertia from the CAD mass properties and from the sizing calculator; confirm the ballast station count on the built wheel. & $I_w \geq$ \qty{6.2352e-4}{\kilogram\metre\squared}; CAD and calculator agree. & \textbf{Verified} (analysis, \qty{100.6}{\percent}) \\
PER-07 & A, T & Spin-up test on the built wheel at end-of-discharge pack voltage (WBS E10). & Design speed reached and sustained on a depleted pack, not only a charged one. & Open \\
PER-08 & A & Derive the achievable ceiling from $K_v$, winding resistance, modulation derate and end-of-discharge bus voltage; compare against the design speed. & Design speed $\leq$ \qty{80}{\percent} of the derived ceiling. & \textbf{Verified} (analysis, \qty{57.8}{\percent}) \\
PER-09 & T & Log driver current throughout balancing and disturbance recovery. & Peak within the \qty{9}{\ampere} continuous rating. & Partial (rig) \\
PER-10 & T & Spin-down test: brake from the design speed with the event logged at a rate resolving a few-millisecond transient. & Arrest within \qty{100}{\milli\second}. & Open \\
PER-11 & T & Record a telemetry session; count logged samples per second of each field. & $\geq$ \qty{50}{\hertz} for state, command and wheel speed. & Open \\
\midrule

\multicolumn{5}{@{}l}{\textbf{Design constraint}}\\
\midrule
CON-01 & I & Measure the built frame against the CAD driving parameter. & Outer edge \qty{150}{\milli\metre} within the stated tolerance. & Open \\
CON-02 & I & Weigh the integrated cube, ready to run. & $\leq$ \qty{1.85}{\kilo\gram}. & Open \\
CON-03 & I & Inspect the assembled machine against the CAD model tree. & Three wheels on mutually orthogonal axes through a common centre. & Open \\
CON-04 & I, D & Inspect the power architecture; run a full demonstration on pack power alone. & No tether or bench supply connected during the demonstration. & Open \\
CON-05 & I & Inspect the bus route and termination against the layout drawings. & Linear chain, short stubs, termination at exactly the two ends. & Open \\
CON-06 & A, T & Sum the coincident worst-case rail load (telemetry burst, servo stall, full sensor load); measure at bring-up. & Within the regulator's derated rating as installed. & Open \\
CON-07 & I & Check the assembled electronics against the frozen mount interface. & Fits the outline, hole pattern and keep-out heights without modification. & Open \\
\end{longtable}
\end{footnotesize}

\paragraph{Evidence pointers.} Verified and partial entries are supported by
Section~\ref{sec:test_1dof} (1-DoF balance duration, recovery angle and driver
current), Section~\ref{sec:massbudget} and Table~\ref{tab:rw_params}
(wheel inertia against target), and the wheel-speed derivation behind
Table~\ref{tab:jumpup_budget}. Section~\ref{sec:test_3dof} carries the result
templates for the cube; every requirement listed Open above resolves through
one of them, and the templates are laid out to mirror the 1-DoF quantities so
the two articles remain directly comparable.
